\documentclass[11pt,a4paper]{article}

\usepackage[utf8]{inputenc}
\usepackage[T1]{fontenc}
\usepackage[margin=2.5cm]{geometry}
\usepackage{RiScVTeX} % !!!

\begin{document}

\begin{center}
    \Huge \RiScVTeX \\ \par \vspace{2mm}
     \normalsize Conde Lucento, 2026 \par \vspace{4mm}
    \large \textit{The Professional RISC-V Documentation Suite}
\end{center}

\listofcodes 
\par
\vspace{4mm}
\hrule
\par
\vspace{4mm}
\risctitle{1. Tipografía y Branding}
El comando principal de identidad es \RiScVTeX. Se puede usar en mitad de una frase y mantiene el espaciado correcto gracias a \texttt{xspace}. 

\subrisctitle{Inline Commands}
Podemos referenciar registros con la función \(\backslash\)reg como \reg{a0} o \reg{t1} manteniendo la coherencia de color. También es posible definir rangos de bits con la función \(\backslash\)bits para manuales ISA: la instrucción \texttt{jal} utiliza el rango de bits \bits{31:12}.

\risctitle{2. Programación en Ensamblador}
A continuación, un ejemplo de código con título, numeración y resaltado de sintaxis:

\begin{riscv}[Rutina de Suma Vectorial]
.data
    vector: .word 1, 2, 3, 4 # Definicion de datos
.text
.globl main
main:
    la   t0, vector     # Carga direccion base
    lw   t1, 0(t0)      # t1 = vector[0]
    addi t1, t1, 10     # Sumar inmediato
    sw   t1, 4(t0)      # Guardar en vector[1]
    
    # Bucle de control
    li   t2, 0          # Contador
loop:
    beq  t2, a0, end    # Si t2 == a0, salir
    addi t2, t2, 1      # Incrementar
    j    loop
end:
    ret
\end{riscv}

\riscwarn{Cuidado con las dependencias de datos (\textit{Data Hazards}) tras un \texttt{lw}. Es recomendable insertar un \texttt{nop} o reordenar el código.}
\newpage 

\risctitle{3. Análisis de Pipeline}
El flujo de una instrucción a través de las etapas del procesador se puede representar así:
\begin{center}
    IF \riscvstep ID \riscvstep EX \riscvstep MEM \riscvstep WB
\end{center}

\risctitle{4. Mapa de Memoria (Stack)}
El entorno \texttt{riscvstack} permite visualizar la pila de forma técnica y limpia:

\begin{riscvstack}[Estado de la Pila en la llamada a Función]
    \riscvmem{0x7ffffffc}{0x00000001}{Variable $x$}%
    \riscvmem{0x7ffffff8}{0x00400560}{Dirección de Retorno \reg{ra}}%
    \riscvmem{0x7ffffff4}{0x7ffffffc}{Frame Pointer \reg{s0} \riscvsp}%
    \riscvmem{0x7ffffff0}{0xDEADBEEF}{Dato basura}%
\end{riscvstack}

\subrisctitle{Notas de implementación}
El paquete \RiScVTeX utiliza \texttt{tabularx} para asegurar que el mapa de memoria ocupe todo el ancho disponible, manteniendo el contenido centrado y resaltado con el color de fondo definido.

\end{document}